\section{Prescribed-Performance Fencing Controller Design}
\label{sec:design}

This section constructs a distributed controller that simultaneously realizes
P4 (enclosure within a prescribed 2-norm envelope), the hard collision fence
P2, and velocity match P3. The design has three building blocks: (i) each
agent estimates the swarm centroid via consensus; (ii) a prescribed-performance
backstepping law is designed for the \emph{centroid} double-integrator
dynamics, using the single scalar transformed error $\xi$ from
\eqref{eq:xi-dyn}; (iii) these terms are assembled into a distributed control
law that superimposes local collision-avoidance and velocity-match terms
without disturbing the centroid control. The controller is summarized in
\eqref{eq:ui}, and its stability/performance guarantees are proved in
\S\ref{sec:analysis}. (Known limitations of the present formulation are
collected honestly in \S\ref{subsec:limits}.)

\subsection{Centroid estimation}
\label{subsec:consensus}

Each agent maintains a local estimate $\hat p_i$ of the swarm centroid
$\bar p=\frac1N\sum_j p_j$ through standard average consensus on positions,
so that $\hat p_i(t)\to \bar p(t)$ (exponentially, or in finite time using
the finite-time centroid enhancer of the companion notes). Agent $i$ then
forms its estimate of the centroid error
\[
\hat{\tilde p}_i = \hat p_i - p_0, \qquad
\dot{\hat{\tilde p}}_i = \dot{\hat p}_i - v_0.
\]
In the centralized description below we write $\tilde p,\dot{\tilde p}$ for
the (estimated) centroid error; in the distributed implementation every
agent evaluates the same expressions from its own $\hat{\tilde p}_i$, which
coincide under consensus.

\subsection{PPC backstepping for the centroid}
\label{subsec:ppc-design}

Apply the transformation of \S\ref{subsec:ppc} to the centroid error norm:
$e=\norm{\tilde p}$, $z=e/\rho$, $\xi=\operatorname{artanh}(z)$,
with the scalar performance function $\rho(t)$ and
$\zeta=1/(1-z^2)>0$, so that
$\dot\xi = \frac{\zeta}{\rho}(\dot e - z\dot\rho)$ and
$\dot e = \frac{\tilde p^T\dot{\tilde p}}{\norm{\tilde p}}$ for
$\tilde p\neq0$.

\paragraph{Step 1 (virtual centroid velocity).}
Treat $\dot{\tilde p}$ as a virtual control. To make $\dot\xi=-c_1\xi$
($c_1>0$) we require the desired norm derivative
\[
\dot e^* = \rho(1-z^2)(-c_1\xi) + z\dot\rho.
\]
Choose the virtual centroid velocity as the vector
\begin{equation}
\alpha_1 = \dot e^*\,\frac{\tilde p}{\norm{\tilde p}}
= \bigl[\rho(1-z^2)(-c_1\xi) + z\dot\rho\bigr]\,
  \frac{\tilde p}{\norm{\tilde p}},
\qquad (\tilde p\neq 0),
\label{eq:alpha1}
\end{equation}
so that $\frac{\tilde p^T}{\norm{\tilde p}}\alpha_1 = \dot e^*$ and hence
$\dot\xi=-c_1\xi$. The apparent singularity at $\tilde p=0$ is removable:
as $\tilde p\to0$, $\dot e^*\to -c_1 e = -c_1\norm{\tilde p}$, so
$\alpha_1\to -c_1\tilde p$, which is bounded; we therefore take
$\alpha_1(0)=0$ by continuous extension. \emph{Note, however, that
$\dot\alpha_1$ (required in \eqref{eq:ubar}) is not differentiable at the
origin; this smoothness caveat is discussed in \S\ref{subsec:limits}.}

\paragraph{Step 2 (centroid control).}
Define the velocity error $s=\dot{\tilde p}-\alpha_1$. Then
$\dot s = \ddot{\tilde p} - \dot\alpha_1$. From
\eqref{eq:centroid-dyn}, with the target acceleration $u_0$ assumed known
(and canceled by feedforward in the distributed law), the centroid dynamics
take the form
\[
\ddot{\tilde p} = \bar u + \delta,
\]
where $\bar u$ is the (yet to be designed) average control
$\bar u=\frac1N\sum_i u_i$, and
$\delta = \frac1N\sum_i d_i - d_0$ is a bounded disturbance satisfying
$\norm{\delta}\le\bar\delta$ (the residual consensus error is folded into
$\bar\delta$). Since
$\dot e = \frac{\tilde p^T}{\norm{\tilde p}}(\alpha_1+s)
       = \dot e^* + \frac{\tilde p^T}{\norm{\tilde p}}s$,
the transformed error obeys
\[
\dot\xi = -c_1\xi + \frac{\zeta}{\rho}\,
          \frac{\tilde p^T}{\norm{\tilde p}}\,s.
\]

Consider $V=\frac12\xi^2+\frac12\norm{s}^2$. Then
$\xi\dot\xi = -c_1\xi^2 + \frac{\zeta}{\rho}\xi\,
\frac{\tilde p^T}{\norm{\tilde p}}s$. Choose
\begin{equation}
\bar u = \dot\alpha_1 - c_2 s
       - \frac{\zeta}{\rho}\,\frac{\tilde p}{\norm{\tilde p}}\,\xi
       - k_r\,\mathrm{sat}(s/\varepsilon),
\qquad c_2>0,
\label{eq:ubar}
\end{equation}
where $\mathrm{sat}(\cdot)$ is the component-wise saturation
($\mathrm{sat}(y)=y$ for $|y|\le1$, $\mathrm{sat}(y)=\mathrm{sign}(y)$
otherwise) and $k_r\ge\bar\delta$, $\varepsilon>0$. Then
\[
s^T\dot s = s^T(\bar u + \delta - \dot\alpha_1)
= -c_2\norm{s}^2
  - \frac{\zeta}{\rho}s^T\frac{\tilde p}{\norm{\tilde p}}\xi
  - k_r\,s^T\mathrm{sat}(s/\varepsilon) + s^T\delta.
\]
The cross terms cancel because
$\frac{\zeta}{\rho}\xi\frac{\tilde p^T}{\norm{\tilde p}}s
 = \frac{\zeta}{\rho}s^T\frac{\tilde p}{\norm{\tilde p}}\xi$, yielding
\[
\dot V = -c_1\xi^2 - c_2\norm{s}^2
        - k_r\,s^T\mathrm{sat}(s/\varepsilon) + s^T\delta.
\]
Since $s^T\delta\le\bar\delta\norm{s}\le k_r\norm{s}$ and
$-k_r\,s^T\mathrm{sat}(s/\varepsilon)\le -k_r\norm{s}+k_r\varepsilon n/2$,
\begin{equation}
\dot V \le -c_1\xi^2 - c_2\norm{s}^2 + \tfrac{n}{2}k_r\varepsilon,
\label{eq:Vdot}
\end{equation}
so $\xi,s$ are uniformly ultimately bounded; taking $\varepsilon\to0$ (or a
signum implementation) recovers asymptotic boundedness, and PPC feasibility
$\rho(0)>\norm{\tilde p(0)}$ guarantees $\xi$ never escapes, i.e.\
$\norm{\tilde p(t)} < \rho(t)$ for all $t\ge0$ (P4).

\subsection{Distributed control law}
\label{subsec:distributed}

The average control $\bar u$ from \eqref{eq:ubar} is realized distributively.
Each agent computes (from its centroid estimate) the common term $\bar u$
and the target feedforward $u_0$, and adds a zero-mean local term:
\begin{equation}
u_i = \bar u + u_0
      - k_p(p_i - \hat p_i) - k_v(v_i - \dot{\hat p}_i)
      + \phi_i + \psi_i,
\label{eq:ui}
\end{equation}
where $\phi_i$ is the collision-avoidance APF of \S\ref{subsec:apf},
$\psi_i$ is a configuration/connectivity term chosen so that
$\sum_i\psi_i=0$ (e.g.\ a zero-sum formation-error term), and
$k_p,k_v>0$. Because $\sum_i\phi_i=0$ and $\sum_i\psi_i=0$, while
$\sum_i(p_i-\hat p_i)\to0$ and $\sum_i(v_i-\dot{\hat p}_i)\to0$ under
consensus, the local terms are zero-sum asymptotically; hence
$\frac1N\sum_i u_i = \bar u + u_0 + o(1)$ and the centroid dynamics reduce
to the analyzed form. The local terms keep each agent near the (estimated)
centroid---yielding the enclosing shape---and, together with
$\dot{\hat p}_i\to\dot{\bar p}\to v_0$, enforce velocity match P3.

\paragraph{Role of each term.}
\begin{itemize}
    \item $\bar u$ (PPC backstepping) $\Rightarrow$ P4 (enclosure within
    $\rho(t)$) and thereby P1 (convex-hull fencing);
    \item $\phi_i$ (APF fence) $\Rightarrow$ P2 (hard collision avoidance);
    \item $-k_p(p_i-\hat p_i)-k_v(v_i-\dot{\hat p}_i)+\psi_i$
    $\Rightarrow$ enclosing formation and P3 (velocity match).
\end{itemize}

\subsection{Prescribed-time variant}
\label{subsec:pt}

To reach the envelope by a user-specified finite time $T$, replace the
constant gain $c_1$ in \eqref{eq:alpha1} by the time-varying gain
\begin{equation}
a(t) = \frac{\pi}{2T}\sec^2\!\left(\frac{\pi t}{2T}\right),
\qquad t\in[0,T),
\label{eq:at}
\end{equation}
and use the finite-time performance function \eqref{eq:rho-pt}. Then
$\alpha_1 = [\rho(1-z^2)(-a(t)\xi)+z\dot\rho]\frac{\tilde p}{\norm{\tilde p}}$,
giving $\dot\xi = -a(t)\xi$ whose solution is
$\xi(t)=\xi(0)\cos(\frac{\pi t}{2T})$; hence $\xi(T)=0$ and
$\xi(t)\equiv0$ for $t\ge T$. In \eqref{eq:ubar} the term
$-c_1\xi$ is replaced by $-a(t)\xi$, and the Lyapunov derivative
\eqref{eq:Vdot} becomes
$\dot V = -a(t)\xi^2 - c_2\norm{s}^2 + \tfrac{n}{2}k_r\varepsilon$ with
$-a(t)\xi^2 = -\frac{\pi}{2T}\xi(0)^2$ bounded as $t\to T^-$.
As with all prescribed-time schemes, the gains $a(t)$ and the corresponding
terms in $\dot\alpha_1$ grow without bound as $t\to T^-$; this must be
respected by actuator limits (or by choosing $T$ conservatively). The
idealised claim of \emph{exact} zero enclosure at $t=T$ is discussed in
\S\ref{subsec:limits}.

\subsection{Stability and Performance Analysis}
\label{sec:analysis}

We first state the standing assumptions and then prove that the closed loop
satisfies P2--P4 (and hence P1).

\begin{assumption}
\label{as:target}
The target acceleration and the agent disturbances are bounded,
$\norm{u_0}\le\bar u_0$, $\norm{d_i}\le\bar d$, $\norm{d_0}\le\bar d_0$,
so that the aggregate centroid disturbance
$\delta=\frac1N\sum_i d_i-d_0$ satisfies $\norm{\delta}\le\bar\delta$.
\end{assumption}

\begin{assumption}
\label{as:feas}
(Initial PPC feasibility) The prescribed performance function satisfies
$\rho(0)>\norm{\tilde p(0)}$.
\end{assumption}

\begin{assumption}
\label{as:consensus}
Each agent obtains a centroid estimate $\hat p_i$ with
$\sum_i(p_i-\hat p_i)\to0$, $\sum_i(v_i-\dot{\hat p}_i)\to0$, and the
configuration term is chosen zero-sum, $\sum_i\psi_i=0$.
\end{assumption}

\begin{theorem}
\label{thm:main}
Under Assumptions~\ref{as:target}--\ref{as:consensus}, the closed-loop system
with the control law \eqref{eq:ui} satisfies, for all $t\ge0$,
\begin{enumerate}
    \item[\textbf{P4}] $\norm{\tilde p(t)}<\rho(t)$ (prescribed 2-norm
    envelope);
    \item[\textbf{P2}] $\norm{p_{ij}(t)}>d$ for all $i\neq j$ (collision
    avoidance);
\end{enumerate}
and, as $t\to\infty$,
\begin{enumerate}
    \item[\textbf{P1}] $\dist(p(t),p_0(t))\to0$ (convex-hull fencing);
    \item[\textbf{P3}] $\norm{v_i(t)-v_0(t)}\to0$ for all $i$ (velocity
    match).
\end{enumerate}
\end{theorem}

\begin{proof}
\noindent\textbf{(i) Prescribed performance P4 and enclosure P1.}
From the backstepping design in \S\ref{subsec:ppc-design}, the Lyapunov
function $V_\xi=\frac12\xi^2+\frac12\norm{s}^2$ satisfies \eqref{eq:Vdot},
i.e.
\[
\dot V_\xi \le -c_1\xi^2 - c_2\norm{s}^2 + \tfrac{n}{2}k_r\varepsilon.
\]
Since $V_\xi\ge\frac12\xi^2$, we have $-c_1\xi^2\ge -2c_1V_\xi$, and dropping
the stabilizing $-c_2\norm{s}^2$ term gives
$\dot V_\xi\le -2c_1V_\xi+\frac{n}{2}k_r\varepsilon$. By the comparison
lemma,
\[
V_\xi(t)\le
\Bigl(V_\xi(0)-\frac{n k_r\varepsilon}{4c_1}\Bigr)e^{-2c_1t}
+ \frac{n k_r\varepsilon}{4c_1},
\]
so $\xi,s$ are exponentially uniformly bounded; in particular $\xi(t)$ can
never diverge. By the feasibility Assumption~\ref{as:feas},
$\rho(0)>\norm{\tilde p(0)}$, hence $z(0)=e(0)/\rho(0)<1$ and
$\xi(0)=\operatorname{artanh}(z(0))$ is finite. Because $\xi(t)$ remains
finite for all $t$, we have $z(t)=\tanh\xi(t)\in(-1,1)$, and therefore
\[
\norm{\tilde p(t)} = \rho(t)\,|z(t)| < \rho(t), \qquad \forall\,t\ge0,
\]
which is P4. Moreover, taking $\varepsilon\to0$ (or a signum
implementation, $k_r\ge\bar\delta$) yields $\xi\to0$ and $s\to0$
exponentially, so $\norm{\tilde p(t)}=\rho(t)\tanh\xi(t)\to0$. Since
$\dist(p,p_0)\le\norm{\tilde p}$ (see \eqref{eq:dist}),
$\dist(p,p_0)\to0$, i.e.\ P1 holds.

\noindent\textbf{(ii) Collision avoidance P2.}
Consider the barrier-type Lyapunov function
\begin{equation}
V_c = \frac12\sum_{i\in\NN}\norm{v_{i0}}^2
    + \frac12\sum_{i\in\NN}\sum_{j\in\NN_i(t)}
      \int_{\norm{p_{ij}}}^\mu \alpha_r(s)\,ds,
\label{eq:Vc}
\end{equation}
where $\alpha_r$ satisfies A1--A2 of \S\ref{subsec:apf}. Its derivative
along the trajectories is
\[
\dot V_c = \sum_{i\in\NN} v_{i0}^T(u_i+d_i-u_0-d_0)
         - \sum_{i\in\NN} v_{i0}^T\phi_i,
\]
where we used the symmetry identity
\(\frac12\sum_i\sum_{j\in\NN_i}
   \alpha_r(\norm{p_{ij}})\frac{p_{ij}^T(v_i-v_j)}{\norm{p_{ij}}}
 = \sum_i v_{i0}^T\phi_i\)
(which follows from the antisymmetry of the pairwise repulsion and
$\sum_i\phi_i=0$). Substituting the control law \eqref{eq:ui} and noting
$u_i-u_0-\phi_i = \bar u - k_p(p_i-\hat p_i)-k_v(v_i-\dot{\hat p}_i)+\psi_i$
(with $\sum_i\psi_i=0$ by Assumption~\ref{as:consensus}) gives
\[
\dot V_c = N\bar u^T\dot{\tilde p}
         - k_p\sum_{i} v_{i0}^T(p_i-\hat p_i)
         - k_v\sum_{i} v_{i0}^T(v_i-\dot{\hat p}_i)
         + \sum_{i} v_{i0}^T(\psi_i+d_i-d_0).
\]
Using $v_i=v_{i0}+v_0$,
\[
-k_v\sum_i v_{i0}^T(v_i-\dot{\hat p}_i)
= -k_v\sum_i\norm{v_{i0}}^2 - k_v\sum_i v_{i0}^Tv_0
  + k_v\sum_i v_{i0}^T\dot{\hat p}_i.
\]
From part (i), $\xi,s$ are bounded, so $\dot{\tilde p}=\alpha_1+s$ is
bounded and, with the designed $\bar u$ of \eqref{eq:ubar}, $\bar u$ is
bounded; by Assumptions~\ref{as:target}--\ref{as:consensus} the remaining
terms $(\psi_i,d_i,p_i-\hat p_i,\dot{\hat p}_i,v_0)$ are bounded. Let
$\norm{v_{i0}}\!:=\!\sqrt{\sum_i\norm{v_{i0}}^2}$. By Cauchy--Schwarz the
bounded terms are majorized by $\beta\norm{v_{i0}}+\gamma$ for some
$\beta,\gamma>0$, while
$-k_v\sum_i v_{i0}^Tv_0\le \frac{k_v}{2}\norm{v_{i0}}^2+\frac{k_vN}{2}\norm{v_0}^2$.
Hence
\[
\begin{aligned}
  \dot V_c
  &\le -k_v\norm{v_{i0}}^2 + \frac{k_v}{2}\norm{v_{i0}}^2
            + \beta\norm{v_{i0}}
            + \Bigl(\gamma+\tfrac{k_vN}{2}\norm{v_0}^2\Bigr)\\
  &= -\frac{k_v}{2}\norm{v_{i0}}^2 + \beta\norm{v_{i0}} + M,
\end{aligned}
\]
for a constant $M>0$. Since the quadratic dominates the linear term,
$\dot V_c\le -c\norm{v_{i0}}^2+M$ for some $c>0$; consequently $V_c$ is
uniformly bounded. In particular the barrier term
$\frac12\sum_i\sum_{j\in\NN_i}\int_{\norm{p_{ij}}}^\mu\alpha_r(s)\,ds$ stays
bounded. By A1,
$\lim_{s\to d^+}\int_s^\mu\alpha_r(\sigma)\,d\sigma=+\infty$, so a bounded
integral forces $\norm{p_{ij}(t)}>d$ for all active pairs and all
$t\ge0$; since pairs entering the safety zone only enlarge the sum, P2 holds
for every pair.

\noindent\textbf{(iii) Velocity match P3.}
From part (i), $\xi\to0$ and $s\to0$, so the centroid velocity error
$\dot{\tilde p}=\alpha_1+s\to0$, i.e.\ $\dot{\bar p}\to v_0$. Define the
local velocity error $\eta_i=v_i-\dot{\hat p}_i$. Using \eqref{eq:ui},
\[
\dot\eta_i = u_i-\ddot{\hat p}_i
= \bar u + u_0 - k_p(p_i-\hat p_i) - k_v\eta_i + \phi_i + \psi_i
  - \ddot{\hat p}_i.
\]
As $t\to\infty$: by Assumption~\ref{as:consensus} $p_i-\hat p_i\to0$ and
$\dot{\hat p}_i\to\dot{\bar p}\to v_0$; by part (i) the centroid control
$\bar u\to0$ (for the exponential performance function the gains
$\dot\rho,\ddot\rho\to0$ and $\xi,s\to0$); and the agents settle into the
enclosing formation so that $\phi_i\to0$ (distances exceed the activation
radius $\mu$) and $\psi_i\to0$. Hence
$\dot\eta_i = -k_v\eta_i + o(1)$, which implies $\eta_i\to0$, i.e.\
$v_i\to\dot{\hat p}_i\to v_0$. Therefore $\norm{v_i-v_0}\to0$ for all $i$,
which is P3.

\noindent\textbf{Prescribed-time case.}
If the performance function \eqref{eq:rho-pt} and the gain $a(t)$ of
\eqref{eq:at} are used, then under the idealised condition $s\equiv0$ and no
disturbance, $\xi(t)=\xi(0)\cos(\frac{\pi t}{2T})$ so that $\xi(T)=0$ and
$\norm{\tilde p(T)}=0<\rho(T)=\bar\rho$; for $t\ge T$ the control may be
switched to the steady formation hold ($\bar u=0$), preserving P2 and P3 and
giving exact enclosure at the user-specified time $T$. The control gains
diverge as $t\to T^-$, which must be respected by actuator limits. (With the
full controller the convergence is only asymptotic/finite-time; see
\S\ref{subsec:limits}.)
\end{proof}

\subsection{Limitations and discussion}
\label{subsec:limits}

The formulation above follows a standard prescribed-performance backstepping
+ artificial-potential-fence architecture and is suitable as a clear
presentation of the idea; it is \emph{not yet} a fully rigorous closed-form
guarantee. The following gaps are acknowledged explicitly.

\begin{itemize}
    \item \textbf{Non-smoothness at the origin.}
    The transformation uses the error \emph{norm} $e=\norm{\tilde p}$ and the
    virtual velocity contains the unit vector $\tilde p/\norm{\tilde p}$,
    which is not differentiable at $\tilde p=0$. While $\alpha_1$ itself
    admits a continuous extension ($\alpha_1\to -c_1\tilde p$), the control
    law \eqref{eq:ubar} requires $\dot\alpha_1$, and $\dot\alpha_1$ carries
    a $1/\norm{\tilde p}$ factor so it is in general unbounded in any
    neighborhood of the origin. The P4 proof therefore implicitly assumes the
    trajectory avoids needing a well-defined $\dot\alpha_1$ at $\tilde p=0$.
    A fully smooth remedy is to transform the \emph{squared} norm
    $e=\tfrac12\norm{\tilde p}^2$ (so $\dot e=\tilde p^T\dot{\tilde p}$ is
    smooth everywhere) and choose $\alpha_1=\dot e^*\,\tilde p/\norm{\tilde p}^2$,
    which is linear and smooth at the origin.

    \item \textbf{Prescribed-time convergence is idealised.}
    The formula $\xi(t)=\xi(0)\cos(\pi t/2T)$ and the claim of enclosure
    ``exactly at $t=T$'' hold only when $\dot\xi=-a(t)\xi$ is satisfied
    exactly, i.e.\ when the velocity-error subsystem is identically zero
    ($s\equiv0$) and there is no disturbance. In the actual backstepping
    controller the extra term $\frac{\zeta}{\rho}\frac{\tilde p^T}{\norm{\tilde p}}s$
    and the robustifying saturation remain, so $\xi$ converges to zero only
    asymptotically (or in finite time under further conditions); the envelope
    $\norm{\tilde p}<\rho(t)$ is still guaranteed for all $t$, but exact zero
    at the prescribed $T$ is not.

    \item \textbf{Consensus--control coupling is not analyzed jointly.}
    The proof treats $\frac1N\sum_i u_i=\bar u+u_0+o(1)$ and folds the
    consensus residual into the disturbance bound $\bar\delta$. A rigorous
    statement would use a composite Lyapunov function for the coupled
    estimation--control dynamics, or explicitly assume the centroid estimate
    converges on a faster time scale than the control loop.

    \item \textbf{Initial collision-free configuration is assumed implicitly.}
    Only the PPC feasibility $\rho(0)>\norm{\tilde p(0)}$ is stated; the
    barrier proof also requires $\norm{p_{ij}(0)}>d$, otherwise the barrier
    integral in $V_c$ is infinite at $t=0$. This should be added as an
    explicit assumption.

    \item \textbf{$\dot\alpha_1$ is not written out.}
    The control law uses $\dot\alpha_1$ without giving its expression. For a
    complete design one should either differentiate $\alpha_1$ explicitly or
    replace the backstepping step by a command filter / dynamic-surface
    control, which avoids computing $\dot\alpha_1$ and also mitigates the
    origin non-smoothness in practice.

    \item \textbf{Robustification introduces a boundary layer.}
    The saturation term $-k_r\,\mathrm{sat}(s/\varepsilon)$ leaves a residual
    $\tfrac n2k_r\varepsilon$ in $\dot V_\xi$; taking $\varepsilon\to0$ (or
    using a signum) removes it but reintroduces chattering. The resulting PPC
    bound is therefore conservative by $O(\varepsilon)$.
\end{itemize}
